\section{Transition Certificates}
\label{sec:method}
In this section, we propose a new type of proof certificates, called transition certificates, to verify that a system $\mathfrak{S}$ satisfies an LTL formula $\phi$. The core idea of transition certificates is to reduce the verification problem to analyzing the behavior of accepting transitions in the NBA $\mathcal{A}$ for $\neg\phi$. By the definition of the transition-based acceptance condition, a word is accepted by $\mathcal{A}$ if and only if accepting transitions occur infinitely often along one of its runs. Therefore, to prove $\mathfrak{S} \models_{\mathfrak{L}} \phi$, it suffices to show that for all trajectories of $\mathfrak{S}$, every accepting transition in $\mathcal{A}$ occurs at most finitely many times. We introduce two types of transition certificates: \emph{transition safety certificates} and \emph{transition persistence certificates}. The two certificates address the problem from fundamentally different angles and concern the set $Acc^{start}$, which comprises the starting states of all accepting transitions in $\mathcal{A}$ (cf. equation~\ref{atstart}).

\begin{itemize}
    \item A \emph{transition safety certificate} proves that the system's trajectories can never reach $Acc^{start}$, thereby guaranteeing that related accepting transitions \textbf{never occur}.
    \item A \emph{transition persistence certificate} allows the system to reach $Acc^{start}$ but proves that related accepting transitions can only be taken \textbf{finitely often}.
\end{itemize}

The high-level intuition is as follows. Given an NBA for $\neg\phi$, we examine each accepting transition individually. If the source state of an accepting transition is unreachable from the initial product state, we certify this with a safety argument, using a function that is non-negative initially and negative at the source state, while preserving non-negativity along trajectories. If the source state is reachable, we take a different approach, showing that the system's dynamics force a strictly positive decrease in a bounded, non-negative function each time the accepting transition could fire, so it can only fire finitely many times. By covering all accepting transitions through one mechanism or the other, we guarantee that no accepting run exists.

\subsection{Transition Safety Certificate}

The objective of a transition safety certificate is to prevent any accepting transition from occurring. To achieve this, we verify that accepting transitions never occur by showing that $Acc^{start}$ is unreachable in the product of the system and the automaton.

\begin{definition}[transition safety certificate]
  \label{transition-safety-certificate-definition}
  Given a system $\mathfrak{S} = (\mathcal{X}, \mathcal{X}_0, f)$, a labelling function $\mathfrak{L}:\mathcal{X} \rightarrow 2^{AP}$ and an NBA $\mathcal{A} = (2^{AP}, Q, q_0, \delta, Acc)$. Construct the product system $\mathfrak{S} \times \mathcal{A} = (\mathcal{Y}, \mathcal{Y}_0, F)$, where $\mathcal{Y} = \{(x, q) \,|\, x \in \mathcal{X}, q \in Q \}$, $\mathcal{Y}_0 = \{ (x_0, q_0) \,|\, x_0 \in \mathcal{X}_0 \}$, and $F(x,q) = \{(x^{\prime}, q') \,|\, x^{\prime} = f(x) \land (q, \mathfrak{L}(x), q') \in \delta\}$. For $q_k \in Acc^{start}$, define the unsafe set $\mathcal{Y}_u = \{(x, q_k) \,|\, x \in \mathcal{X}\}$. A bounded function $B_{k}(x, q): \mathcal{X} \times Q \rightarrow \mathbb{R}$ is a \emph{transition safety certificate} with respect to $q_k \in Acc^{start}$ if for all $(x,q) \in \mathcal{Y}$, $(x_0, q_0) \in \mathcal{Y}_0$, $(x^{\prime}, q^{\prime}) \in F(x, q)$, we have:
  \begin{align}
    &B_{k}(x_0, q_0) \geq 0, \label{btsc1} \\
    &B_{k}(x, q_k) < 0, \label{btsc2} \\
    &B_{k}(x, q) \geq 0 \Rightarrow B_{k}(x^{\prime}, q^{\prime}) \geq 0. \label{btsc3}
  \end{align}
\end{definition}

% Intuitively, these conditions ensure that: 1) the certificate value is non-negative at the initial state; 2) it becomes negative precisely on the unsafe set $\mathcal{Y}_u^k$ associated with $q_k$; and 3) once non-negative, its value remains non-negative along any trajectory of the product system. The certificate $B_k$ guarantees $q_k$ is unreachable, preventing all accepting transitions from $q_k$. The following lemma formalizes this guarantee. Its proof is provided in the Appendix~\ref{appendix-proof}.

Intuitively, condition (\ref{btsc1}) guarantees that the certificate value is non-negative at the initial state of the product system; condition (\ref{btsc2}) ensures that the certificate becomes negative precisely on the unsafe set $\mathcal{Y}_u^k$ associated with $q_k$, thereby marking states where accepting transitions could originate; condition (\ref{btsc3}) implies that once the certificate value is non-negative at a state, it remains non-negative along any trajectory of the product system, thus preserving the safety invariant.

Collectively, these conditions ensure that the certificate $B_k$ guarantees $q_k$ is unreachable, preventing all accepting transitions from $q_k$. The following lemma formalizes this guarantee. Its proof is provided in Appendix~\ref{appendix-proof}.

\begin{lemma}
  \label{transition-safety-lemma}
  If a transition safety certificate $B_k$ with respect to $q_k \in Acc^{start}$ can be found, then all runs of words in $\mathfrak{L}(\mathfrak{S})$ never visit state $q_k$, thereby preventing all accepting transitions starting from $q_k$.
\end{lemma}

The next theorem follows directly from applying the above lemma to every state in $Acc^{start}$. If a transition safety certificate can be found for each state in $Acc^{start}$, then no accepting transition can be taken, ensuring the LTL specification holds.

\begin{theorem}
  \label{transition-safety-LTL}
  Given a system $\mathfrak{S} = (\mathcal{X}, \mathcal{X}_0, f)$, a labelling function $\mathfrak{L}:\mathcal{X} \rightarrow 2^{AP}$ and an NBA $\mathcal{A} = (2^{AP}, Q, q_0, \delta, Acc)$ for $\neg\phi$. If for any $q_k \in Acc^{start}$, a transition safety certificate $B_k$ can be found, then it can be guaranteed that $\mathfrak{S} \models_{\mathfrak{L}} \phi$.
\end{theorem}

\begin{proof}
By leveraging Lemma~\ref{transition-safety-lemma}, all start states of accepting transitions are unreachable. Consequently, all accepting transitions never occur, ensuring that the system's language is not accepted and thereby guaranteeing $\mathfrak{S} \models_{\mathfrak{L}} \phi$.
\end{proof}

\subsection{Transition Persistence Certificate}
In contrast, the \emph{transition persistence certificate} offers an alternative approach. Instead of completely preventing the system from reaching states where accepting transitions may happen, this certificate proves that even if such states are reached, the corresponding accepting transitions can occur only a finite number of times. This is achieved by requiring a decrease in a certificate value each time an accepting transition is taken.

\begin{definition}[transition persistence certificate]
  \label{transition-persistence-certificate-definition}
  Given a system $\mathfrak{S} = (\mathcal{X}, \mathcal{X}_0, f)$, a labelling function $\mathfrak{L}:\mathcal{X} \rightarrow 2^{AP}$ and an NBA $\mathcal{A} = (2^{AP}, Q, q_0, \delta, Acc)$. For $(k, l) \in Acc^{pair}$, define the indicator function $I_{k,l}(x) = 1$ if $(q_k, \mathfrak{L}(x), q_l) \in Acc^{k,l}$, and $0$ otherwise. A bounded function $B_{k, l}(x): \mathcal{X} \rightarrow \mathbb{R}$ is a \emph{transition persistence certificate} with respect to $Acc^{k,l}$ if there exists $\epsilon > 0$ such that for all $x_0 \in \mathcal{X}_0$, $x \in \mathcal{X}$, $x^{\prime} = f(x)$, we have:
  \begin{align}
    &B_{k, l}(x_0) \geq 0, \label{btpc1} \\
    &B_{k, l}(x) \geq 0 \Rightarrow B_{k, l}(x^{\prime}) \geq 0, \label{btpc2} \\
    &B_{k, l}(x) \geq B_{k, l}(x^{\prime}) + I_{k,l}(x)\epsilon. \label{btpc3}
  \end{align}
\end{definition}

We emphasize two important aspects of Definition~\ref{transition-persistence-certificate-definition}. First, each certificate $B_{k,l}$ is constructed for a \emph{specific} accepting transition pair $(k, l) \in Acc^{pair}$; different pairs may require different certificates. Second, $B_{k,l}(x): \mathcal{X} \rightarrow \mathbb{R}$ is a function of the system state $x$ alone and does not depend on the current automaton state. This design is sound because the indicator $I_{k,l}(x)$ captures a \emph{necessary condition} for the accepting transition $(q_k, \mathfrak{L}(x), q_l)$ to occur. Regardless of which automaton state a particular run occupies at time step $i$, whenever $I_{k,l}(x_i) = 0$ the transition from $q_k$ to $q_l$ cannot be taken. By bounding the number of visits to the enabling region $\{x \mid I_{k,l}(x) = 1\}$ along system trajectories, the certificate provides an upper bound on the number of times the accepting transition can actually occur in \emph{any} run.

This design introduces a degree of conservatism. Because $B_{k,l}$ constrains the total number of time steps during which $I_{k,l}(x) = 1$ without regard to the current automaton state, it limits visits to the enabling region even when the system occupies a state other than $q_k$. As a consequence, a valid persistence certificate may fail to exist even when the LTL specification is satisfied, if the trajectory frequently enters the enabling region while the automaton is in a state other than $q_k$. This conservatism is the price of working with a function over $\mathcal{X}$ rather than the full product space $\mathcal{X} \times Q$, and it is analogous to the over-approximation inherent in trajectory-based Lyapunov arguments.

Intuitively, condition (\ref{btpc1}) guarantees that the certificate $B_{k,l}$ is non-negative at the initial state of the system; condition (\ref{btpc2}) ensures that $B_{k,l}$ remains non-negative along any trajectory of the system; condition (\ref{btpc3}) implies that $B_{k,l}$ decreases by at least $\epsilon$ whenever the state's label matches an accepting transition from $q_k$ to $q_l$, and remains non-increasing otherwise.

% By (\ref{btpc1}) and (\ref{btpc2}), $B_{k,l}$ remains non-negative along any trajectory starting from $x_0 \in \mathcal{X}_0$. By (\ref{btpc3}), $B_{k,l}$ decreases by at least $\epsilon$ each time $I_{k,l}(x) = 1$. 

Since accepting transitions from $q_k$ to $q_l$ can only occur when $(q_k, \mathfrak{L}(x), q_l) \in Acc^{k,l}$, the number of such transitions is bounded by $B_{k,l}(x_0)/\epsilon$, ensuring finite occurrence. This yields the following lemma. Its proof is provided in the Appendix~\ref{appendix-proof}.

\begin{lemma}
  \label{transition-persistence-lemma}
  If a transition persistence certificate $B_{k,l}$ can be found, then all runs of words in $\mathfrak{L}(\mathfrak{S})$ make accepting transitions from $q_k$ to $q_l$ happen finitely often.
\end{lemma}

The next theorem is obtained by applying the lemma to every accepting transition pair.

\begin{theorem}
  \label{transition-persistence-LTL}
  Given a system $\mathfrak{S} = (\mathcal{X}, \mathcal{X}_0, f)$, a labelling function $\mathfrak{L}:\mathcal{X} \rightarrow 2^{AP}$ and an NBA $\mathcal{A} = (2^{AP}, Q, q_0, \delta, Acc)$ for $\neg\phi$. if for any $(k, l) \in Acc^{pair}$, a corresponding transition persistence certificate $B_{k,l}$ can be found, then it can be guaranteed that $\mathfrak{S} \models_{\mathfrak{L}} \phi$.
\end{theorem}

\begin{proof}
By leveraging Lemma~\ref{transition-persistence-lemma}, all accepting transitions occur only finitely often, which ensures that the system's language is not accepted, thereby guaranteeing $\mathfrak{S} \models_{\mathfrak{L}} \phi$.
\end{proof}

While both the methods established by Theorem~\ref{transition-safety-LTL} and Theorem~\ref{transition-persistence-LTL} can be used to verify LTL properties, the two types of certificates are complementary and, when combined, significantly expand the range of systems that can be verified. In some scenarios, requiring the system to completely avoid certain states might be too strict, making transition safety certificates impossible to find. However, the system might still satisfy the property if the accepting transitions occur only finitely often, a condition that can be certified by a transition persistence certificate. Conversely, consider verifying $\mathbf{G} \neg a$ using the NBA for $\mathbf{F} a$, where $Acc^{0,0} = \{(q_0, \emptyset, q_0), (q_0, \{a\}, q_0)\}$. For the persistence certificate $B_{0,0}$, condition (\ref{btpc3}) requires $B_{0,0}(x) \geq B_{0,0}(f(x)) + \epsilon$ for any $x$ where an accepting transition is enabled. This would demand an infinite descent in the bounded certificate value, which is impossible, making it impossible to satisfy the conditions for a transition persistence certificate, and thus such a certificate cannot exist. Yet, a transition safety certificate can successfully verify the property by proving that state $q_0$ (the start of the accepting transitions) is unreachable in the product system. Thus, providing both certificates covers a wider range of system behaviors, enhancing the method's applicability. We now present the main theorem for LTL verification. For each accepting transition start state in the NBA, we use either a transition safety certificate or transition persistence certificates to prove accepting transitions from that state occur finitely often.

\begin{theorem}[LTL verification]
  \label{ltl-verification-theorem}
  Given a system $\mathfrak{S} = (\mathcal{X}, \mathcal{X}_0, f)$, a labelling function $\mathfrak{L}: \mathcal{X} \rightarrow \Sigma$ and an NBA $\mathcal{A} = (\Sigma, Q, q_0, \delta, Acc)$ representing the negation of an LTL formula $\phi$. For each $q_k \in Acc^{start}$, if we can find either:
  \begin{enumerate}
    \item A transition safety certificate $B_k$ with respect to $q_k$, or
    \item Transition persistence certificates $B_{k,l}$ for all $(k, l) \in Acc^{pair}$ where $q_k$ is the start state,
  \end{enumerate}
  then it can be guaranteed that $\mathfrak{S} \models_{\mathfrak{L}} \phi$.
\end{theorem}

\begin{proof}
  Given an accepting transition start state $q_k \in Acc^{start}$:
  \begin{enumerate}
    \item If a transition safety certificate $B_k$ is found, by Lemma~(\ref{transition-safety-lemma}), all runs never visit state $q_k$, thereby preventing all accepting transitions starting from $q_k$.
    \item For each $(k, l) \in Acc^{pair}$ where the start state is $q_k$, if a transition persistence certificate $B_{k,l}$ is found, by Lemma~(\ref{transition-persistence-lemma}), all runs make accepting transitions from $q_k$ to $q_l$ happen finitely often.
  \end{enumerate}
  This ensures that all accepting transitions in $Acc$ occur only finitely often. By the definition of NBA acceptance, a word $w$ is accepted if and only if there exists a run where accepting transitions occur infinitely often, i.e., $\mathit{Inf}(r) \cap Acc \neq \emptyset$. Since all accepting transitions occur only finitely often, for any word $w \in \mathfrak{L}(\mathfrak{S})$ and any run $r$ on $w$, we have $\mathit{Inf}(r) \cap Acc = \emptyset$. Therefore, no word in $\mathfrak{L}(\mathfrak{S})$ is accepted by automaton $\mathcal{A}$, implying $\mathfrak{L}(\mathfrak{S}) \cap \mathfrak{L}(\mathcal{A}) = \emptyset$. Since $\mathfrak{L}(\mathcal{A}) = \mathfrak{L}(\neg \phi)$, we have $\mathfrak{L}(\mathfrak{S}) \subseteq \mathfrak{L}(\phi)$, and consequently $\mathfrak{S} \models_{\mathfrak{L}} \phi$.
\end{proof}

One structural aspect of Theorem~\ref{ltl-verification-theorem} warrants clarification. For each $q_k \in Acc^{start}$, the verification must commit to a single strategy: either supply a transition safety certificate $B_k$ that simultaneously prevents all accepting transitions originating from $q_k$, or supply individual persistence certificates $B_{k,l}$ for every pair $(k, l) \in Acc^{pair}$ with start state $q_k$. The two strategies cannot be mixed for the same $q_k$, since a safety certificate $B_k$ makes a global argument that $q_k$ is entirely unreachable, which already prevents all accepting transitions from $q_k$ and leaves no residual role for a persistence certificate at that state.

In summary, the combination of transition safety certificates and transition persistence certificates offers a more flexible verification framework. We emphasize that the novelty of our approach does not lie in merely combining barrier-like and ranking-like arguments, which are well-established proof techniques. Rather, the key contribution is the \emph{transition-level decomposition} of the verification problem. Instead of constructing a single monolithic certificate over the entire product state space, we decompose the task at the granularity of individual accepting transitions. For each accepting transition start state $q_k$, we independently select the most suitable proof strategy (safety or persistence) based on the local structure of the automaton and the system dynamics. This per-transition decomposition yields certificates that each operate over a smaller domain with simpler conditions, reducing both the dimensionality of the synthesis problem and the restrictiveness of the sufficient conditions. The resulting framework maintains a smaller search space compared to closure certificates while providing two distinct mechanisms to ensure that accepting transitions occur only finitely often, ultimately proving that the system satisfies its LTL specification.

